\documentclass{article} \usepackage[utf8]{inputenc} \usepackage{fancyhdr} \usepackage{verbatim} \usepackage[brazil]{babel} \usepackage[T1]{fontenc} \title{Relatório de Análise de Portas e Serviços com Nmap e Tcpdump} \author{Samuel Abrao} \date{\today} \begin{document} \maketitle \section{Metodologia do Ambiente e Ferramentas} Para este laboratório, a arquitetura de duas máquinas virtuais (VMs) foi replicada utilizando contêineres Docker, o que permitiu uma simulação precisa do ambiente sem a necessidade de recursos computacionais adicionais de VMs completas. A escolha pelo Docker foi estratégica, pois ele oferece isolamento de serviços de forma leve, resolvendo problemas comuns de conflitos de portas que poderiam ocorrer ao rodar serviços como Apache e MySQL diretamente no sistema operacional local. A aplicação web (Apache/PHP) foi hospedada em um contêiner, e o banco de dados (MySQL) em outro, ambos comunicando-se através de uma rede interna gerenciada pelo Docker. Essa abordagem espelhou a arquitetura de "VM Cliente e VM Servidor" proposta no laboratório. Para a análise de rede, o utilitário \texttt{tcpdump} foi utilizado em vez do Wireshark. O tráfego entre os contêineres Docker trafega em uma rede virtual interna, que não é acessível diretamente pelas interfaces de rede do sistema hospedeiro. Portanto, o \texttt{tcpdump} foi instalado e executado dentro do próprio contêiner do servidor web, o que permitiu a captura direta dos pacotes de dados enviados para o banco de dados. Essa metodologia garantiu que a prova da comunicação e da falha de serviço pudesse ser obtida de forma confiável e precisa, superando a incompatibilidade do Wireshark com essa arquitetura de rede. \section{Análise e Resultados} \subsection{Análise de Portas e Serviços com Nmap} O primeiro conjunto de exercícios teve como objetivo familiarizar-se com os comandos básicos do Nmap para identificar portas e serviços em um alvo de teste, \texttt{scanme.nmap.org} (IP: 45.33.32.156). \subsubsection{Verificação de Sinal de Vida do Host} O comando \texttt{nmap -sn} foi utilizado para realizar um escaneamento de ping, que confirma se o host está online. \begin{verbatim} samuelabrao@SamuelDeb:~/VSCODE/IDP/4 Semestre/IDP-Redes-2025.2/A1/lab03$ nmap -sn 45.33.32.156 Starting Nmap 7.93 ( https://nmap.org ) at 2025-08-27 17:50 -03 Nmap scan report for scanme.nmap.org (45.33.32.156) Host is up (0.21s latency). Nmap done: 1 IP address (1 host up) scanned in 0.35 seconds \end{verbatim} O resultado confirma que o host está online e acessível, permitindo a continuidade dos testes. \subsubsection{Scan Básico de Portas} Um escaneamento básico foi realizado nas 1000 portas TCP mais comuns para identificar serviços abertos. \begin{verbatim} samuelabrao@SamuelDeb:~/VSCODE/IDP/4 Semestre/IDP-Redes-2025.2/A1/lab03$ nmap 45.33.32.156 Starting Nmap 7.93 ( https://nmap.org ) at 2025-08-27 17:52 -03 Nmap scan report for scanme.nmap.org (45.33.32.156) Host is up (0.21s latency). Not shown: 993 closed tcp ports (conn-refused) PORT STATE SERVICE 22/tcp open ssh 80/tcp open http 135/tcp filtered msrpc 139/tcp filtered netbios-ssn 445/tcp filtered microsoft-ds 9929/tcp open nping-echo 31337/tcp open Elite Nmap done: 1 IP address (1 host up) scanned in 22.38 seconds \end{verbatim} O escaneamento revelou que o host possui as portas 22 (SSH), 80 (HTTP), 9929 e 31337 abertas, enquanto outras portas estão filtradas, indicando a presença de um firewall. \subsubsection{Verificação de Portas Específicas} Para uma análise mais focada, as portas 22, 80 e 443 foram escaneadas especificamente. \begin{verbatim} samuelabrao@SamuelDeb:~/VSCODE/IDP/4 Semestre/IDP-Redes-2025.2/A1/lab03$ nmap -p 22,80,443 45.33.32.156 Starting Nmap 7.93 ( https://nmap.org ) at 2025-08-27 17:50 -03 Nmap scan report for scanme.nmap.org (45.33.32.156) Host is up (0.31s latency). PORT STATE SERVICE 22/tcp open ssh 80/tcp open http 443/tcp closed https Nmap done: 1 IP address (1 host up) scanned in 0.65 seconds \end{verbatim} O resultado confirmou que as portas 22 e 80 estão abertas e a porta 443 está fechada. \subsubsection{Detecção de Versões de Serviços} O escaneamento de detecção de versão (\texttt{nmap -sV}) foi utilizado para identificar as versões dos serviços rodando nas portas abertas. \begin{verbatim} samuelabrao@SamuelDeb:~/VSCODE/IDP/4 Semestre/IDP-Redes-2025.2/A1/lab03$ nmap -sV 45.33.32.156 Starting Nmap 7.93 ( https://nmap.org ) at 2025-08-27 17:54 -03 Nmap scan report for scanme.nmap.org (45.33.32.156) Host is up (0.35s latency). Not shown: 979 closed tcp ports (conn-refused) PORT STATE SERVICE VERSION 22/tcp open ssh OpenSSH 6.6.1p1 Ubuntu 2ubuntu2.13 (Ubuntu Linux; protocol 2.0) ... (resultados omitidos para concisão) 9929/tcp open nping-echo Nping echo 31337/tcp open tcpwrapped Service Info: OS: Linux; CPE: cpe:/o:linux:linux_kernel Nmap done: 1 IP address (1 host up) scanned in 133.31 seconds \end{verbatim} Este escaneamento revelou que a porta 22 está executando o OpenSSH em um sistema Linux Ubuntu, fornecendo informações valiosas sobre o ambiente do host. \subsection{Monitoramento com Tcpdump} O \texttt{tcpdump} foi utilizado para monitorar o tráfego da porta 3306, provando a conectividade com o banco de dados e evidenciando a falha de comunicação quando o serviço foi interrompido. \subsubsection{Cenário de Sucesso} O log abaixo demonstra um cadastro bem-sucedido. A sequência de pacotes mostra o \textit{handshake} de três vias do TCP (pacotes SYN, SYN-ACK e ACK) seguido pela troca de dados (\texttt{P} - Push) e, por fim, o encerramento da conexão (\texttt{F} - FIN). \begin{verbatim} root@e805c72ae38c:/var/www/html# tcpdump -i any 'tcp port 3306' tcpdump: WARNING: any: That device doesn't support promiscuous mode (Promiscuous mode not supported on the "any" device) tcpdump: verbose output suppressed, use -v[v]... for full protocol decode listening on any, link-type LINUX_SLL2 (Linux cooked v2), snapshot length 262144 bytes 01:05:11.824661 eth0 Out IP e805c72ae38c.58876 > lab3-mysql.app_default.3306: Flags [S] ... length 0 01:05:11.824691 eth0 In IP lab3-mysql.app_default.3306 > e805c72ae38c.58876: Flags [S.] ... length 0 01:05:11.824697 eth0 Out IP e805c72ae38c.58876 > lab3-mysql.app_default.3306: Flags [.] ... length 0 ... (pacotes de dados e finalização) 01:05:11.832716 eth0 Out IP e805c72ae38c.58876 > lab3-mysql.app_default.3306: Flags [F.] ... length 0 01:05:11.832725 eth0 In IP lab3-mysql.app_default.3306 > e805c72ae38c.58876: Flags [F.] ... length 0 01:05:11.832730 eth0 Out IP e805c72ae38c.58876 > lab3-mysql.app_default.3306: Flags [.] ... length 0 \end{verbatim} \subsubsection{Cenário de Falha} No cenário de erro, com o servidor MySQL parado, o \texttt{tcpdump} não registrou pacotes de resposta. Isso ocorreu porque o pacote de solicitação de conexão foi enviado, mas não havia um serviço escutando na outra extremidade para responder. A ausência de pacotes de resposta é a evidência clara de que o serviço do banco de dados estava indisponível, impedindo que o \textit{handshake} de conexão fosse concluído. \section{Conclusão} O laboratório demonstrou a eficácia do Nmap e do Tcpdump para a análise e monitoramento de redes. O Nmap permitiu uma visão detalhada dos serviços e portas do host alvo, revelando sua configuração de segurança e a presença de firewalls. O Tcpdump, por sua vez, foi essencial para a análise de pacotes em um ambiente de contêineres, provando a conectividade do sistema web com o banco de dados e evidenciando a falha de comunicação quando o serviço foi interrompido. Esses resultados confirmam a realização e o entendimento dos objetivos propostos no laboratório. \end{document}